\documentclass[12pt,a4paper]{article}

\usepackage[czech]{babel}
\usepackage[T1]{fontenc}
\usepackage[utf8]{inputenc}

\usepackage{geometry}
\usepackage{twemojis}
\usepackage[most]{tcolorbox}

\geometry{margin=2.5cm}

\title{\textbf{Rozbor díla „451 stupňů Fahrenheita“}}
\author{}
\date{}

\begin{document}
\maketitle

\section*{Autor}

\begin{tcolorbox}[breakable]

\textbf{Autor:}
\begin{itemize}
\item Ray Bradbury
\end{itemize}

\textbf{Název díla:}
\begin{itemize}
\item 451 stupňů Fahrenheita
\end{itemize}

\textbf{Zařazení autora do uměleckého směru, charakteristika směru:}
\begin{itemize}
\item Poválečná americká literatura 20. století, sci-fi literatura a antiutopie. 
Díla často varují před nebezpečím totalitní společnosti, manipulací lidí a zneužíváním technologií.
\end{itemize}

\textbf{Určení století, v němž autor tvořil:}
\begin{itemize}
\item 20. století
\end{itemize}

\textbf{Určení dalších autorů stejného směru + dílo:}
\begin{itemize}
\item George Orwell – 1984
\item Aldous Huxley – Konec civilizace (Brave New World)
\end{itemize}

\textbf{Další autorova díla:}
\begin{itemize}
\item Marťanská kronika
\item Pampeliškové víno
\item Zlatá jablka slunce
\end{itemize}

\textbf{Specifický styl autora:}
\begin{itemize}
\item Autor píše především sci-fi literaturu. Jeho styl je obrazný a atmosférický. 
Často kritizuje společnost, která se vzdává přemýšlení a nahrazuje ho zábavou a technologiemi.
\end{itemize}

\end{tcolorbox}

\section*{Charakteristika uměleckého textu jako celku}

\begin{tcolorbox}[breakable]

\textbf{Literární druh, žánr a forma:}
\begin{itemize}
\item epika, próza, román (sci-fi, antiutopie)
\end{itemize}

\textbf{Typ vypravěče:}
\begin{itemize}
\item er-forma, personální vypravěč
\end{itemize}

\textbf{Dominantní slohový postup:}
\begin{itemize}
\item vyprávěcí
\end{itemize}

\textbf{Vysvětlení názvu díla:}
\begin{itemize}
\item Název označuje teplotu, při které se samovolně vznítí papír (451 °F). 
V románu mají hasiči za úkol pálit knihy, proto název symbolizuje ničení literatury a vědění.
\end{itemize}

\textbf{Posouzení aktuálnosti díla:}
\begin{itemize}
\item Dílo je aktuální i dnes. Ukazuje společnost zahlcenou zábavou a rychlými informacemi, 
která ztrácí schopnost přemýšlet. Některé prvky lze přirovnat k dnešním médiím a sociálním sítím.
\end{itemize}

\textbf{Určení místa a času děje:}
\begin{itemize}
\item neurčitá budoucnost, pravděpodobně ve Spojených státech amerických
\end{itemize}

\textbf{Stručné nastínění děje:}
\begin{itemize}
\item Hlavní hrdina Guy Montag pracuje jako hasič, jehož úkolem je pálit knihy. 
Postupně začíná pochybovat o systému, který zakazuje čtení. 
Setkání s dívkou Clarissou v něm probudí zájem o skutečný život a přemýšlení. 
Montag začne knihy tajně číst a hledá smysl života mimo manipulovanou společnost.
\end{itemize}

\textbf{Smysl díla:}
\begin{itemize}
\item Kritika společnosti, která zakazuje přemýšlení a nahrazuje knihy zábavou a povrchními informacemi. 
Dílo varuje před manipulací společnosti a ztrátou svobody myšlení.
\end{itemize}

\textbf{Pravděpodobný adresát:}
\begin{itemize}
\item široká veřejnost, především moderní společnost ohrožená manipulací médií
\end{itemize}

\textbf{Kompozice:}
\begin{itemize}
\item chronologická
\end{itemize}

\break

\textbf{Zařazení do kontextu autorovy tvorby:}
\begin{itemize}
\item Jedno z nejznámějších děl autora. Typické sci-fi dílo, které kritizuje společnost a varuje před nebezpečím technologií.
\end{itemize}

\textbf{Porovnání s filmovou verzí:}
\begin{itemize}
\item Existuje filmová adaptace z roku 1966. Některé scény jsou oproti knize změněny, například setkání Montaga s Clarissou. 
Film se liší i v některých detailech děje.
\end{itemize}

\end{tcolorbox}

\section*{Úryvek k rozboru}
Hlavní je klid, Montagu. Vypisuj pro
lidi soutěže, ve kterých dostanou cenu za to, že se naučí zpaměti populární šlágr, nebo hlavní města
různých států, nebo kolik tun obilí sklidili loni v Iowě. Napěchuj do nich údaje, které nejsou
nebezpečné, zahlť je tolika "fakty", až se budou cítit přecpaní, ale budou mít dojem, že doslova srší
informacemi. Potom se jim bude zdát, že myslí. Budou mít pocit pohybu, aniž se pohnou, a budou
spokojení, protože fakta tohoto druhu se nemění. Jenom jim nedávat nic tak ošemetného jako
filozofii nebo sociologii, aby si uměli dát dohromady dvě a dvě. To je cesta k nespokojenosti.
Každý, kdo dovede rozebrat a zase složit telestěnu — a dnes to dokáže většina lidí — je šťastnější
než člověk, který se pokouší s logaritmickým pravítkem vypočítat, změřit nebo dát do rovnic
vesmír. Ten nemůžeš změřit nebo dát do rovnice, pokud nechceš, aby se člověk cítil jako zvíře a
úplně osamocený. Já to znám, zkoušel jsem to. K čertu s tím!

\begin{tcolorbox}[breakable]

\textbf{Atmosféra úryvku:}
\begin{itemize}
\item napjatá, manipulativní
\end{itemize}

\textbf{Počet postav:}
\begin{itemize}
\item Guy Montag
\item kapitán Beatty
\end{itemize}

\textbf{Charakteristika postav:}

\textbf{Guy Montag}
\begin{itemize}
\item hlavní hrdina románu
\item pracuje jako hasič, který pálí knihy
\item postupně začíná pochybovat o systému a hledá smysl života
\end{itemize}

\textbf{Kapitán Beatty}
\begin{itemize}
\item velitel hasičů
\item inteligentní a vzdělaný člověk
\item podporuje režim a snaží se Montaga přesvědčit, že knihy jsou nebezpečné
\end{itemize}

\textbf{Další postavy}

\textbf{Clarissa McClellanová}
\begin{itemize}
\item mladá dívka, která Montaga přiměje přemýšlet o světě
\item je zvídavá a citlivá
\end{itemize}

\textbf{Mildred Montagová}
\begin{itemize}
\item manželka Montaga
\item tráví většinu času sledováním telestěn a poslechem mušlí
\item představuje typického člověka této společnosti
\end{itemize}

\textbf{Profesor Faber}
\begin{itemize}
\item bývalý profesor
\item pomáhá Montagovi pochopit význam knih
\end{itemize}

\textbf{Vztahy mezi postavami:}
\begin{itemize}
\item Beatty podezřívá Montaga z toho, že čte knihy
\item vztah mezi nimi je napjatý a nepřátelský
\end{itemize}

\textbf{Zařazení úryvku do kontextu díla:}
\begin{itemize}
\item Úryvek pochází z první části knihy, kdy kapitán Beatty navštíví Montaga a vysvětluje mu principy fungování společnosti.
\end{itemize}

\textbf{Použité jazykové prostředky:}
\begin{itemize}
\item metafory
\item opakování
\item ironie
\end{itemize}

\end{tcolorbox}
\section*{Děj}
\subsection*{Část první: Ohniště a salamandr}

Román vypráví o hlavní postavě jménem \textbf{Guy Montag}, který pracuje jako požárník. V tomto světě však požárníci nehasí požáry, ale naopak \textbf{pálí knihy}. Montag nosí na uniformě číslo \textbf{451} a znak \textbf{salamandra}. Je ženatý s \textbf{Mildred Montagovou}, která tráví většinu času sledováním televizních stěn a posloucháním hudby pomocí malých sluchátek, kterým se říká \textit{mušle}. Mildred sní o tom, že si koupí čtvrtou televizní stěnu, aby byla ještě více ponořená do světa zábavy. \\
\\
Jednou cestou domů Montag potká zvláštní šestnáctiletou dívku jménem \textbf{Clarissa}. Má bledou pleť, světlé vlasy a modré oči. Montagovi připadá velmi zvláštní, protože se chová jinak než ostatní lidé. Klade mu jednoduché otázky o jeho práci požárníka a vypráví mu o svém strýci, který jí vyprávěl o tom, jak vypadal život v minulosti. \\
\\
Na požární stanici Montag popisuje \textbf{mechanického ohaře}. Je to osminohý kovový stroj připomínající psa, který dokáže pomocí chemického čichu vyhledávat nepřátele systému. Na stanici Montag mluví s velitelem \textbf{Beattym} a se svými kolegy. \\
\\
Jednoho dne vyjedou k požáru do domu ženy, která ukrývá knihy. Žena však odmítne dům opustit a raději \textbf{zemře spolu se svými knihami}. Požárníci proto zapálí celý dům i s ní uvnitř. Během zásahu Montag tajně \textbf{schová jednu knihu do kabátu}.\\
\\
Po návratu domů začne Montag prožívat vnitřní pochybnosti a psychické napětí. Onemocní horečkou a začne přemýšlet o smyslu své práce. Druhý den zjistí, že už má dvě hodiny zpoždění do práce. Do jeho domu přijde velitel Beatty, který začne vysvětlovat historii požárníků a důvody, proč společnost začala pálit knihy. Montag má velký strach, protože má jednu knihu schovanou pod polštářem. \\
 \\
Po Beattyho odchodu Montag přizná Mildred pravdu. Ukáže jí, že doma tajně schovává knihy. Vysype na zem asi dvacet knih různých barev. Mildred z toho dostane paniku a chce knihy okamžitě spálit. Montag ji však zastaví a řekne jí, že teď jsou oba vinni. \\
 \\
Najednou někdo zaklepe na dveře. Oba však zůstanou potichu a návštěvu ignorují. Nakonec začnou společně číst knihu od \textbf{Jonathana Swifta}. \\

\subsection*{Část druhá: Písek a síto}

Montag a Mildred pokračují ve čtení knih, ale většině z nich nerozumějí. Montag si vzpomene na starého profesora jménem \textbf{Faber}, kterého kdysi potkal v parku. Tehdy Faber rychle schoval knihu do kabátu, protože se bál. Nakonec Montagovi dal své telefonní číslo. \\
 \\
Když Mildred oznámí, že k nim přijdou na návštěvu její kamarádky, Montag se rozhodne odejít z domu a navštívit Fabera. Cestou ve vlaku začne číst \textbf{Bibli}. Snaží se soustředit, ale z reproduktorů neustále zní hlasitá reklama na \textit{Denhamovu zubní pastu}. Montag se nakonec rozkřičí, protože se nedokáže soustředit, a ostatní cestující na něj začnou zírat. \\
 \\
Montag přijde k Faberovi domů a zaklepe na dveře. Faber ho nejprve podezřívá, ale nakonec ho pustí dovnitř. Montag mu ukáže Bibli a požádá ho, aby ho naučil chápat knihy. Faber vysvětluje, že hodnotu knih tvoří tři věci: \textbf{kvalita informací, čas na přemýšlení a možnost tyto myšlenky uplatnit v životě}. \\
 \\
Montag navrhne plán: chtěl by tajně \textbf{tisknout nové knihy}. Faber se nejdříve bojí a odmítá, ale když Montag začne trhat stránky z Bible, nakonec souhlasí. \\
 \\
Faber mu také ukáže malé zelené zařízení, které vypadá jako mušle do ucha. Ve skutečnosti je to \textbf{malé rádio}, pomocí kterého spolu mohou tajně komunikovat na dálku. \\
 \\
Večer přijdou k Montagovým domů tři Mildrediny kamarádky. Baví se o mužích, válce, potratech a politice. Rozhovor je velmi povrchní a bezcitný. Montag se nakonec rozčílí a začne jim číst básně z knihy. Jedna z žen se rozpláče. Aby Montag nevzbudil podezření, Faber mu do ucha radí, aby situaci obrátil v žert. \\
 \\
Později Montag odejde na požární stanici. Vrátí knihu Beattymu, který ji bez zájmu hodí do odpadkového koše. Požárníci začnou hrát poker. Během hry Beatty cituje různé slavné autory a básníky, aby Montaga zmátl. Montag je ve velkém stresu. \\
 \\
Najednou zazní sirény oznamující nový požár. Požárníci nasednou do vozu Salamandr a vyrazí na místo zásahu. Když dorazí, Montag s hrůzou zjistí, že stojí \textbf{před jeho vlastním domem}.

\subsection*{Část třetí: Oheň se rozhoří}

Velitel Beatty přikáže Montagovi, aby \textbf{sám spálil svůj vlastní dům}. Mildred rychle uteče z domu a odjede taxíkem pryč. Montag začíná tušit, že právě ona nahlásila knihy.\\
\\
Montag zapálí dům plamenometem. Beatty mu řekne, že po skončení zásahu půjde Montag do vězení. V návalu vzteku a zoufalství Montag náhle otočí plamenomet proti Beattymu a \textbf{spálí ho zaživa}.\\
\\
Mechanický ohař na něj okamžitě zaútočí. Montag ho sice zničí, ale ohař mu ještě stihne vstříknout jed do nohy. Montag si vzpomene, že na zahradě má ještě jednu knihu, vezme ji s sebou a uteče.\\
\\
Pronásledují ho policie i vrtulníky. Montag se snaží dostat k Faberovi. Když dorazí k jeho domu, vidí v televizi přímý přenos svého pronásledování. Faber mu dá staré oblečení a poradí mu, aby utekl k řece a pokračoval lesem podél staré železniční tratě.\\
\\
Montag skočí do řeky a nechá se unášet proudem. Poprvé po dlouhé době si všimne hvězd na obloze a ticha přírody.\\
\\
Nakonec v lese narazí na skupinu starších mužů. Jsou to bývalí profesoři a intelektuálové, kteří si \textbf{pamatují celé knihy zpaměti}, aby je zachovali pro budoucnost.\\
\\
Krátce poté vypukne válka a město je během několika minut zničeno bombardováním. Skupina lidí se rozhodne vrátit k ruinám města, aby pomohla vybudovat novou společnost založenou na \textbf{vědomostech, paměti a knihách}.



\section*{Názor na dílo}

Román na mě působil velmi silně. Připomíná mi dílo 1984 od George Orwella, protože také zobrazuje totalitní společnost a manipulaci lidí. Kniha se čte dobře a hlavní myšlenka je velmi aktuální i v dnešní době.

\end{document}
